Within the two proposed configurations, the Stage-B threshold estimators were compared separately at Fp1 and at Fp2 — the operating point Experiment~1 identified as consistently best-performing on both corpora (Table~\ref{tab:channel_selection}) — rather than a best-channel-per-session oracle, with each electrode's precision, recall and $F_1$ scored inside the full 32-channel montage run at the 30\,s reference epoch (Table~\ref{tab:exp3_estimator}). At both electrodes and on both corpora, the direction of the metric changes did not support the stated precision hypothesis.
On Internal, Proposed-Med achieved precision/recall/$F_1$ of 86.45\%/88.93\%/86.32\% at Fp1 and 86.28\%/89.74\%/86.82\% at Fp2, compared with 87.04\%/87.43\%/85.55\% (Fp1) and 86.75\%/88.20\%/85.99\% (Fp2) for Proposed-Mean: $F_1$ rose by $+0.76$ (Fp1) and $+0.83$ (Fp2) percentage points, in both cases because a recall gain ($+1.50$ and $+1.54$ points) outweighed a precision loss ($-0.59$ and $-0.48$ points).
On Cao2018, the same pattern held: Proposed-Med's $F_1$ rose by $+0.29$ (Fp1) and $+0.43$ (Fp2) percentage points, with recall gains of $+1.96$ and $+2.14$ points against precision losses of $-1.29$ and $-1.28$ points. Pooled performance followed suit, with $F_1$ increasing by $+0.50$ (Fp1) and $+0.61$ (Fp2) percentage points.
None of the six dataset/electrode combinations reached significance under a two-tailed Wilcoxon signed-rank test on session-level $F_1$ (Internal: $p=0.407$ at Fp1, $p=0.435$ at Fp2; Cao2018: $p=0.331$ at Fp1, $p=0.795$ at Fp2; Pooled: $p=0.791$ at Fp1, $p=0.775$ at Fp2), so, as with the earlier best-channel-per-session comparison, the two estimators remain statistically indistinguishable at either fixed electrode. Pooled, Proposed-Med exceeded Proposed-Mean in 44/104 (Fp1) and 50/104 (Fp2) sessions, was equal in 2 and 4, respectively, and was lower in 58 and 50, indicating a per-session outcome close to even at both electrodes.
